\documentclass[11pt]{article}
\usepackage[a4paper,margin=1in]{geometry}
\usepackage{amsmath}
\usepackage{xcolor}

% ---- a visually distinct block for the referee's comments ----
\definecolor{commentgray}{gray}{0.35}
\newenvironment{refcomment}%
  {\begin{quote}\color{commentgray}\itshape}%
  {\end{quote}}
\newcommand{\response}{\medskip\noindent\textbf{Response.}\ }
\newcommand{\point}[1]{\bigskip\noindent\textbf{#1}}

\setlength{\parskip}{0.4em}
\setlength{\parindent}{0pt}

\title{Response to Referee 1\\[2pt]
\large Manuscript JFM-2026-1781\\
\emph{Strain-coupled one-dimensional turbulence for rapid distortion}\\
Sharma \& Kerstein}
\date{}

\begin{document}
\maketitle

We thank the referee for a careful reading and for identifying the two
directions in which the original manuscript was underdeveloped. The paper
has been substantially rewritten. It is now built around the single-line
analysis the referee identified as the strongest part; an independent
strained-box DNS with exact-RDT companions has been added as the spectral
benchmark of record; the leading-edge application is carried through to an
absolute prediction and compared with four experiments; and the closure
deficiencies the referee named are now measured against the DNS and stated
as quantified limits rather than left open. Section, figure, table and
equation numbers below refer to the revised manuscript.

\point{Major point 1 (validation of the main result outside the model).}
\begin{refcomment}
``The central claim is that at $Sk_t/\varepsilon \approx 0.8$ the spectrum
develops a broadband distortion that differs from the rigid spectral
translation predicted by rapid-distortion theory. This is established by
comparing the model with itself, once with the strain switched on and once
with it switched off. \ldots\ the predicted distorted spectrum is never
directly compared against an independent spectral benchmark.''
\end{refcomment}

\response
We accept the objection. The strain-on/strain-off comparison is retained
only as an internal decomposition (now \S\,4.3, figure~9) and is no longer
offered as validation. The validation of record is now an independent,
projection-matched spectral benchmark: a $128^3$ strained-box DNS of
plane-strained homogeneous turbulence in a deforming-frame pseudo-spectral
formulation, together with exact-RDT companion calculations, all reduced to
the same one-dimensional line observables (\S\,4.6; DNS protocol in
Appendix~B.3). The three-way comparison of the strain-induced spectral
anisotropy---SC-ODT against exact projected RDT against DNS---is
figure~11 and table~1 at $e=1$, $Sk_t/\varepsilon=0.8$. It shows directly
what the model reproduces (the upwash amplification) and where it departs
(its wavenumber-uniform signature against the wavenumber-dependent DNS and
RDT). Section~4.6.2 and figure~12 add a distance-from-linear-theory
envelope (the rms log-residual of a two-parameter exact-RDT-distorted
von~K\'arm\'an fit) as a function of strain rapidity $Sk_t/\varepsilon$,
giving the departure a measured, monotone size. A five-way version at three
rapidities ($Sk_t/\varepsilon=0.4,0.8,16$) is figure~16.

\point{Major point 2 (the original single-line development is not used in
the reported calculations).}
\begin{refcomment}
``The simulations are performed using the closure introduced in Section
3.2, whereas the more complete single-line formulation developed later is
not used in generating the principal results. \ldots\ it is unclear whether
the main spectral conclusions would remain unchanged under the formulation
that the paper ultimately argues is the more appropriate one.''
\end{refcomment}

\response
The referee is correct that the reported spectral calculations use the
algebraic (LRR) moment closure, and we now state this plainly at each
figure and in Section~5.7. Section~5 has been reframed so that it no longer
claims a closed single-line functional that then goes unused. Instead it
(i) states the obstruction exactly---the line carries the
transverse-integrated spectra of equation~(5.3), and the rapid term
requires the direction-resolved projection weight that those integrals
destroy (\S\,5.2); (ii) derives sharp bounds on the rapid term from line
data alone by linear programming, with the residual indeterminacy tabulated
(\S\,5.4, table~2, Appendix~A.4); and (iii) shows that the collapsed
spectral form of equations~(5.6)--(5.9) reduces exactly to the moment
closure in the isotropic limit, equation~(5.11), which is what justifies
using that closure in the simulations (\S\,5.7, figure~14). The consistency
question the referee raises---whether the conclusions survive where the
moment closure loses accuracy---is answered by measurement rather than
assumption: the axisymmetry hypothesis on which the collapse rests is
measured against the DNS (\S\,5.5, table~3, figure~13) and the
component-ordering error is localised (\S\,5.6; next point), while the
envelope of \S\,4.6.2 bounds the resulting distance from linear theory over
the strain range. Exercising the full wavenumber-dependent kernel of
equation~(5.6) inside ODT is identified as the missing ingredient (\S\,6.4)
and left as future work; we do not claim it here.

\point{Major point 3 (the rapid pressure--strain closure becomes
qualitatively inaccurate under plane strain).}
\begin{refcomment}
``Under plane strain, the exact solution predicts growth of the neutral
spanwise component, which eventually overtakes the upwash component,
whereas both closures predict the opposite trend. \ldots\ What remains
unclear is how sensitive the subsequent spectral predictions are to this
known deficiency.''
\end{refcomment}

\response
This is now measured, localised and bounded rather than only acknowledged.
The exact rapid term evaluated in the DNS (\S\,5.6) gives
$\Pi^{(\mathrm r)}_{33}/k_tS=0.14$ at $Sk_t/\varepsilon=0.8$ and $0.27$ at
$16$, the latter exceeding $\Pi^{(\mathrm r)}_{11}=0.23$: the overtaking is
real and sets in at rapid strain. We show that the error lives in the
moment closure's anisotropy-dependent term---LRR reproduces
$\Pi^{(\mathrm r)}_{33}$ at $Sk_t/\varepsilon=0.8$ but delivers half the
exact value at $16$, and IP gives zero identically---and not in the line
model: the ODT moment-level $b_{ij}$ follows the closure to $\sim 2\%$. In
the spectral representation the ordering is carried by the anisotropy
scalar $c$ through the kernels~(5.7)--(5.9), which the moment reduction
discards. The same deficiency appears at the moment level in the
plane-strain $b_{11}$ (\S\,3.4, figure~6): the fast-strain DNS tracks exact
RDT while the closures give $b_{11}\approx0$. On sensitivity: the
acoustically relevant quantity is the upwash amplification $b_{22}$, and
the A1 control experiment (figure~13\textit{b}, table~3) shows it
indistinguishable between plane and axisymmetric strain even at
$Sk_t/\varepsilon=16$, i.e.\ robust to the azimuthal structure that carries
the ordering error; the distance envelope of \S\,4.6.2 gives the overall
sensitivity of the spectral prediction its measured bound.

\point{Major point 4 (the leading-edge connection is not carried through).}
\begin{refcomment}
``First, it is not demonstrated that the spectrum produced by the ODT line
is the same quantity required by Amiet's theory, which is naturally
formulated in terms of chordwise and spanwise wavenumbers. Second, the
strain history responsible for generating the turbulence distortion is
deferred to future work.''
\end{refcomment}

\response
The connection is now carried through in a dedicated Section~7. The
two-dimensional upwash spectrum $\Phi_{ww}(K_x,K_y)$ that Amiet's theory
consumes is supplied from the axisymmetric representation of \S\,5.3 fitted
to the line spectra, so the model's one-dimensional output is mapped to the
quantity the aeroacoustic kernel requires; the computed distorted inflow is
then fed to Amiet's response both as a ratio from which aerofoil geometry
and observer position cancel (figure~17) and as an absolute far-field
prediction at the reference operating point (figure~18). The assembled
chain is compared with four published experiments (figure~19) and lies
within $1$--$2$\,dB rms where the facilities are cleanest; the term omitted
by the frozen-spectrum assumption is $5$--$20$\,dB (figures~17--21). On the
second part, we agree the strain history from a real stagnation flow is not
modelled, and we now scope the claim as a frozen-strain calculation: the
operating point is entered through an effective strain estimated
geometrically (potential-flow stagnation-line accumulation truncated at the
eddy scale, varied over $\Lambda/4$--$\Lambda$; figure~19 caption), and the
distance envelope of \S\,4.6.2 is carried through as the uncertainty
statement (figures~20--21). The mapping from model operating point to the
stagnation environment is stated as the remaining step, not claimed as
solved.

\bigskip
\hrule
\bigskip
\noindent\textbf{Minor points.}

\point{(i) Novelty of pre-impact distortion; dos Santos, Ribeiro, Piccolo.}
\response
The introduction has been rewritten. The pre-impact distortion is now
presented as established, with references to dos Santos \textit{et al.}
(2023), Ribeiro \textit{et al.} (2023) and Piccolo \textit{et al.} (2024),
and the distinction our contribution rests on---prescribing an RDT- or
measurement-derived spectrum versus \emph{evolving} one, including the
nonlinear relaxation---is stated explicitly (\S\,1).

\point{(ii) The dissipation estimate as the residual of two larger terms.}
\response
The residual-of-two-large-terms estimate has been replaced. Because the
eddy events conserve energy and the production is exact, molecular
dissipation is the only sink and is fixed by the resolved energy budget of
equation~(4.1), with every term available on the line (\S\,4.4). This
estimator-free budget gives $\chi=Sk_t/\varepsilon\approx1.2$, corroborated
within a few tens of per cent by the direct gradient estimate
$\langle 2\nu(\partial u_i/\partial y)^2\rangle\approx1.1$; the spectral
peak migration corroborates the $\chi=O(1)$ reading independently of the
dissipation estimate altogether.

\point{(iii) The Lee \& Reynolds strain-parameter identification (factor of
two).}
\response
This is now discussed where the comparison is made (\S\,3.4). Lee \&
Reynolds characterise each run by $S^{*}=Sq^2/\varepsilon$, which equals
$2\chi$ \emph{identically} because $q^2=2k_t$; the manuscript now states
that this is a correspondence of definitions, not a calibration, and notes
that the axisymmetric anisotropy they report is essentially independent of
$S^{*}$, so the comparison target is parameter-free.

\point{(iv) Propositions 1 and 2.}
\response
Agreed. They are demoted from results to consistency checks and presented
as such in the verification section (\S\,3.1--3.2).

\point{(v) Length; Section 2 non-original; implementation detail to
supplementary material.}
\response
The standard-ODT material of \S\,2.1 is compressed to a review keyed to
citations of the original ODT papers, and established-versus-new material
is separated through Section~2. Implementation detail, the CMK comparison
(\S\,4.5 summary; full account in supplementary \S\,S1) and the full
intervention catalogue (supplementary \S\,S2; labels in Appendix~B.1) have
been moved out of the main text.

\bigskip
\noindent
We are grateful for the recommendation to build the study around the
single-line analysis; that is now the organising spine of the revised
paper.

\end{document}
